\section{Related Work}
\label{sec:related_work}

\subsection{LLM Inference Systems}

LLM inference systems have improved performance by changing how requests are scheduled, how model state is stored, and where computation runs. Orca and Sarathi-Serve change the scheduling of request execution~\cite{yu2022orca,agrawal2024sarathi}; vLLM and CacheGen address the management and transfer of KV-cache state~\cite{kwon2023pagedattention,liu2024cachegen}; and SGLang provides a runtime for structured generation~\cite{zheng2024sglang}. Other systems distribute execution or state across devices and memory tiers through offloading, model parallelism, phase separation, migration, or preemption~\cite{sheng2023flexgen,aminabadi2022deepspeed,zhong2024distserve,patel2024splitwise,sun2024llumnix,miao2024spotserve}. QServe illustrates the further demands of low-precision execution on hardware-specific code~\cite{lin2025qserve}. These mechanisms meet in the engines we study: request control, cached state, and hardware execution must work together. Systems papers explain how to build and optimize these mechanisms; our focus is the code changes developers make when engine behavior is wrong.

\subsection{Bugs in AI Infrastructure}

Empirical studies of DL applications and production jobs identify common symptoms and causes across model-development workflows~\cite{zhang2018tensorflowbugs,zhang2020dljobfailures,islam2019bugs,humbatova2020taxonomy}. Work on frameworks, libraries, and production platforms examines faults in reusable infrastructure; Jia et al., for example, also analyze repairs made inside TensorFlow~\cite{yang2022frameworkbugs,chen2023frameworkbugs,tambon2024silentbugs,gao2023platformquality,jia2021tensorflowbugs}. Compiler studies characterize failures in another specialized part of the stack~\cite{shen2021compilerbugs,du2021compilerbugs,shen2025compilerevolution}. These studies establish useful ways to examine symptoms, causes, and fixes, but their subjects do not isolate LLM serving, where concurrent batching, KV-cache state, model loading, APIs, and heterogeneous backends must be maintained together.

The closest study is Liu et al.'s analysis of 929 issue-reported bugs from five inference engines~\cite{liu2026firstlook}. It relates symptoms to root causes, compares bug characteristics across engines, derives fix strategies from associated code changes, measures fix effort using duration and changed lines, and examines how bug types evolve. These are substantial findings about inference-engine reliability. Our unit of analysis is a merged, manually verified bug-fixing PR with its complete patch. This makes it possible to ask different code-level questions: which files and functions repeatedly participate in fixes, how repair operations vary with the architectural layer containing the fault, and whether complete fixes touch more files over time. The larger corpus supports these comparisons across six engines, but the difference is in the maintenance questions being answered, not just the number of cases.

\begin{table*}[t]
  \centering
  \caption{Comparison with empirical studies of bugs in AI infrastructure. Corpus gives analyzed artifacts / subject systems using each study's unit. Sym. = explicitly analyzed symptoms or observable impacts; Loc. = root causes and/or internal fault locations; S--L = explicit symptom/impact--cause/location associations; Rep. = recurring repair operations derived from code changes; Conc. = file/function-level concentration; Scope = patch scope covering both changed-line magnitude and modified-file breadth; Evol. = changes in bug or repair characteristics over calendar/project time. Loc. represents architectural locations, whereas Evol. represents their change over time. A checkmark denotes explicit coverage.}
  \label{tab:empirical_study_comparison}
  \small  
  \setlength{\tabcolsep}{10pt}
  \renewcommand{\arraystretch}{1.12}
  \begin{tabular}{@{}l l r c c c c c c c@{}}
    \toprule
    Study & Domain & Corpus & Sym. & Loc. & S--L & Rep. & Conc. & Scope & Evol. \\
    \midrule
    Islam et al.~\cite{islam2019bugs} & DL applications & 3,216 / 5 & \checkmark & \checkmark & -- & -- & -- & -- & \checkmark \\
    Humbatova et al.~\cite{humbatova2020taxonomy} & DL systems & 1,059 / 3 & -- & \checkmark & -- & -- & -- & -- & -- \\
    Jia et al.~\cite{jia2021tensorflowbugs} & DL library & 202 / 1 & \checkmark & \checkmark & \checkmark & \checkmark & -- & -- & -- \\
    Shen et al.~\cite{shen2021compilerbugs} & DL compilers & 603 / 3 & \checkmark & \checkmark & \checkmark & -- & -- & -- & -- \\
    Du et al.~\cite{du2021compilerbugs} & DL compilers & 2,717 / 5 & \checkmark & \checkmark & \checkmark & -- & -- & -- & \checkmark \\
    Yang et al.~\cite{yang2022frameworkbugs} & DL frameworks & 1,127 / 8 & \checkmark & \checkmark & -- & \checkmark & -- & -- & -- \\
    Chen et al.~\cite{chen2023frameworkbugs} & DL frameworks & 1,000 / 4 & \checkmark & \checkmark & \checkmark & -- & -- & -- & -- \\
    Liu et al.~\cite{liu2026firstlook} & LLM inference engines & 929 / 5 & \checkmark & \checkmark & \checkmark & \checkmark & -- & -- & \checkmark \\
    \midrule
    \textbf{Our work} & \textbf{LLM inference engines} & \textbf{12,707 / 6} & \checkmark & \checkmark & \checkmark & \checkmark & \checkmark & \checkmark & \checkmark \\
    \bottomrule
  \end{tabular}
\end{table*}


Table~\ref{tab:empirical_study_comparison} shows both the overlap and the distinction. Liu et al.~\cite{liu2026firstlook} already relate symptoms to root causes and group fix strategies by those causes. We instead locate faults in architectural layers and examine repair operations within each layer. Liu et al. also analyze changed lines; our patch-scope analysis additionally counts modified files in complete PRs and tracks their breadth over time. Alongside the file- and function-level concentration analysis, this connects observed failures to where repair work recurs and how broadly completed repairs extend through the codebase.

\subsection{Repair and Reliability Techniques}

General software-engineering research has identified recurring fix patterns and patch actions~\cite{pan2009bugfixpatterns,sobreira2018dissection,zhong2015realbugfixes}. Studies of vulnerability concentration, fault location, change history, code churn, and change complexity also show why the distribution and scope of changes matter for quality assurance~\cite{liu2020vulnerabilitydistribution,ostrand2005faultlocation,graves2000changehistory,nagappan2005codechurn,hassan2009changecomplexity,kamei2013jitqa}. These ideas motivate our repair-operation categories and our file-, function-, and patch-level measures. We use them to describe verified repairs in inference engines, rather than to train or evaluate a fault predictor.

Testing research addresses specific ways to expose or diagnose defects. Cross-backend comparison detects inconsistent behavior~\cite{pham2019cradle}; generated models, mined API calls, and fuzzing exercise DL libraries and compilers~\cite{wang2020lemon,wei2022freefuzz,gu2022muffin,li2023comet,guo2020audee,liu2023nnsmith,deng2023titanfuzz}. Other work targets CUDA synchronization, GPU numerical behavior, and silent inference errors~\cite{wu2020simulee,rathnasuriya2025gpunumerical,gu2026ekka}. These methods study how to find particular failures. Our retrospective analysis studies the fixes that were actually merged: the symptoms they addressed, the code regions involved, and the operations and file changes needed to repair them. This evidence can help identify where future testing and repair methods need to focus.

The results also define concrete targets for evaluating such techniques. Backend checks can compare behavior across devices, whereas Scheduler/Runtime checks need request-state traces and regression tests that follow multi-file changes. Our study does not compare these methods; it identifies which code surfaces and repair boundaries their evaluations should cover.
 
